\section{The conical-region error}
\label{app:cone}

This appendix records an error made during the development of this work, why
the verification then in place could not detect it, and how it was eventually
found. It is included because the failure is instructive rather than
idiosyncratic: the construction that produced it is a natural one, the mesh it
produces is valid by every ordinary measure, and the resulting scientific
conclusion was not merely inaccurate but inverted.

\subsection{The construction}
\label{app:cone:construction}

The task is to build a closed mesh bounding the material beyond a plane. A
natural approach, and the one first taken, avoids clipping altogether. Select
the facets whose centroids satisfy $x > c$; these form an open surface patch.
Find the boundary edges of the patch---those belonging to exactly one selected
facet---and close the surface by adding, for each such edge $(i,j)$, the
triangle $(O, \mathbf{v}_j, \mathbf{v}_i)$, where $O$ is a point inside the
body, in our case the mean of the vertices.

The appeal of this is that it requires no geometric predicates beyond a
centroid test, no intersection arithmetic, and no tolerance. The resulting mesh
is closed, consistently oriented, and its volume, computed by the divergence
theorem, agrees with the sum of the signed tetrahedra
$\{(O, \mathbf{v}_a, \mathbf{v}_b, \mathbf{v}_c)\}$ over the selected facets to
machine precision. All of this is true. None of it is relevant.

\subsection{What the mesh actually bounds}
\label{app:cone:solid}

The solid enclosed is not the material beyond the plane. It is the radial cone
with apex at $O$ subtended by the selected surface patch: the union of all
segments joining $O$ to a point of the patch. For a sphere with $O$ at the
centre this is a spherical sector. The two solids share their outer boundary
and nothing else. The plane-cut lobe is bounded below by the plane $x = c$; the
sector is bounded below by a lateral surface running from the rim of the patch
to the interior point $O$, and therefore includes material at every $x$ between
$x_O$ and $c$ that the intended region excludes.

For a sphere of radius $R$ cut at $x = c$, with $O$ at the centre, both solids
admit closed forms, given as Eqs.~\eqref{eq:appcap} and~\eqref{eq:appcone} of
Appendix~\ref{app:clip}. At $R = 1$, $c = 0.5$:
\begin{equation}
\frac{V_{\mathrm{sector}}}{V_{\mathrm{cap}}}
= \frac{2\pi R^{2}h/3}{\pi h^{2}(3R-h)/3}
= \frac{2R^{2}}{h(3R-h)}
= 1.600 ,
\qquad
\bar{x}_{\mathrm{cap}} - \bar{x}_{\mathrm{sector}} = 0.1125\,R .
\label{eq:appconeratio}
\end{equation}
The sector has sixty per cent more volume and its centroid lies more than a
tenth of a radius further in. Figure~\ref{fig:cone} shows the two.

The discrepancy is not a small-parameter effect that vanishes for a thin cap.
As $h \to 0$ the ratio in Eq.~\eqref{eq:appconeratio} diverges as
$2R/(3h)$: the thinner the intended region, the worse the substitution. This
matters directly, because Section~\ref{sec:itokawa:sweep} advances the plane
precisely into that regime.



\subsection{The consequence on Itokawa}
\label{app:cone:consequence}

The density model assigns an elevated density to the region and reduces the
complement to conserve the total mass. With the correct region, raising
$\denshead$ adds mass in the head, displacing the centre of mass toward $+x$.
With the sector, it adds mass along a wedge whose centroid lies at $+0.5625R$
rather than $+0.6750R$ on the sphere analogue, and on Itokawa reaches back past
the neck into the body.

The recovered contrast inverted. Table~\ref{tab:appcone} gives both models as
they were computed with the conical construction, alongside the plane-cut
result at the same plane on the same shape model.

\begin{table}[htbp]
\centering
\small
\caption{The Itokawa inversion computed with the conical region, against the
plane-cut result. The conical runs used a mesh decimated to $41{,}369$ facets
and a body volume of $0.0177848$~km$^{3}$; the plane-cut column is the
reference run of Section~\ref{sec:itokawa:reference} on the unreduced mesh.
Both conical models return a head several times \emph{less} dense than the
body, with the minimum pinned to the lower edge of the density scan.}
\label{tab:appcone}
\begin{tabular}{lrrr}
\toprule
 & \multicolumn{2}{c}{Conical region} & Plane-cut \\
\cmidrule(lr){2-3}\cmidrule(lr){4-4}
Quantity & head & neck & head \\
\midrule
Region facets            & $6{,}412$ & $5{,}730$ & $7{,}950$ \\
Region volume (km$^3$)   & $0.004229$ & $0.001486$ & $0.002922$ \\
Volume fraction          & $23.8\%$ & $8.4\%$ & $16.4\%$ \\
Volume closure error     & $2.1\times10^{-14}\%$ & $0\%$ & $2.1\times10^{-14}\%$ \\
$\densR$ (kg\,m$^{-3}$)  & $300$ & $300$ & $2800$ \\
$\densrest$ (kg\,m$^{-3}$) & $2547$ & $2169$ & $1858$ \\
Contrast                 & $0.118$ & $0.138$ & $1.507$ \\
Interior minimum         & no & no & yes \\
Improvement              & $6.121\%$ & $3.831\%$ & $16.889\%$ \\
\bottomrule
\end{tabular}
\end{table}

Three features of the conical columns are consistent with a region that reaches
into the body. The recovered density sits at $300$~kg\,m$^{-3}$, the lower
limit of the scan, in both models: the objective function is minimized by
removing as much mass as the scan permits from a wedge that spans the neck, and
would have gone lower had it been allowed to. Neither minimum is interior, so
the edge-minimum guard of Section~\ref{sec:failures:edge} would flag both---but
that guard was added later, and at the time the values were reported as
results. And the head region holds $23.8\%$ of the volume against $16.4\%$ for
the plane-cut region at the same plane, the excess being precisely the wedge
between $O$ and the plane.

The volume closure error is the point of the table. At
$2.1\times10^{-14}$ per cent it is as small as floating-point arithmetic allows,
and it is identical for the conical and the plane-cut construction. A check
that returns the same value for a correct solid and for one wrong by a factor
of $1.6$ in volume is measuring something other than correctness.




\subsection{Why the verification did not catch it}
\label{app:cone:verification}

The construction was accompanied by a check described as a volume-closure test.
It compared the volume of the assembled mesh, computed by the divergence
theorem over its facets, against the sum of the signed tetrahedra used to build
it. On the Itokawa head region it agreed to $2\times10^{-14}$ per cent.

The agreement was exact and meaningless. Both quantities describe the same
solid---the sector---and the test verified that the mesh had been assembled
consistently from the tetrahedra, which it had. A test constructed this way
cannot fail for the reason that mattered, and therefore carries no information
about it. Nor would the other checks of Section~\ref{sec:methods:verify} have helped, had
they existed at the time. The complement test is passed by both constructions,
since sectors over complementary patches partition the body just as plane-cut
lobes do. The manifold test is passed by the sector, which is a perfectly good
closed surface. And the containment test finds nothing, because the sector is a
single component and there is no other component for any of its facets to lie
inside. Of the five checks the paper now applies, exactly one distinguishes the
two solids.

What distinguishes them is a statement about their vertices rather than their
volumes, their topology or their embedding. Every vertex of the plane-cut region
satisfies $x \ge c$ by construction. The sector contains $O$, which does not.
Testing
\begin{equation}
\max_i \,(c - x_i) \le \epsilon
\label{eq:apppredicate}
\end{equation}
over the vertices of the region returns $0$ for the lobe and $0.5$ for the
sector in the units of the sphere test, where the plane sits at $x = 0.5$ and
$O$ at the origin. The test costs three lines and inspects the property that
was actually in doubt.

\subsection{How it was found}
\label{app:cone:discovery}

Not by any check. The conical results were internally consistent, passed every
test then applied, and gave a smooth, well-behaved dispersion curve. They were
questioned because the recovered contrast disagreed in \emph{sign} with
\citet{lowry2014} and \citet{kanamaru2019}, both of which report a denser head.

The diagnosis followed from constructing the same region on a unit sphere,
where the closed forms of Appendix~\ref{app:clip:analytic} are available. The
tetrahedral construction returned a volume of $1.047453$ and a centroid of
$0.562213$, matching the sector formulae to five and four decimals
respectively, and matching the spherical cap not at all.

Three observations follow that are worth carrying beyond this particular error.

A disagreement with published work is not a verification method, but it is
sometimes the only signal available. Had Itokawa's true structure been unknown,
or had the sign happened to come out right, the error would have survived.

Analytic test cases are cheap and decisive. A sphere has closed-form caps,
slabs and sectors; the entire diagnosis took one script and no shape model. The
same lesson recurs in Section~\ref{sec:discussion:ellipsoid}, where an analytic
ellipsoid settles in one experiment a question that mesh-to-mesh comparison had
answered wrongly. It is worth constructing such a case before trusting a
numerical routine, not after.

And a verification test should be designed against a specific hypothesis of
failure. The question is not whether a check passes, but what it would look
like if the code were wrong in a way one has not thought of. The volume-closure
test answered the question ``was the mesh assembled consistently?''; nobody had
asked ``is this the solid I meant?'' The same pattern recurred later in the
work, with the buried facets of Section~\ref{sec:failures:buried}: the manifold
test, added precisely because volume comparison had proved insufficient, was
itself insufficient against a failure nobody had yet imagined. Each check
answers the question its author posed, and the list of questions is never
known to be complete.

\subsection{Status of the superseded results}
\label{app:cone:status}

All results computed with the conical construction are superseded, and none
appear in this paper except in Table~\ref{tab:appcone}, where they are shown
for comparison. They are retained in the archived results file
(Section~\ref{sec:conclusions}, data availability) under a key marking them as
superseded, together with the diagnosis, so that the record of the work is
complete and the error can be reproduced by anyone wishing to check this
account.

One further point should be recorded. The pilot study preceding this work,
whose results informed Paper~IV, was internally inconsistent in a related way:
it evaluated the potential on a planar-capped head mesh while computing the
region volume as a sum of origin-anchored tetrahedra, which is a sector volume.
The mismatch is the likely origin of a three per cent volume discrepancy noted
at the time and left unexplained. No conclusion of Paper~IV depends on it,
since that paper reports no density inversion, but the episode is recorded here
for completeness.


